% ============================================================
%  Parting Is a Sweet Sorrow
%  An Ambient Co-Presence Space for Friendships Drifting Apart
%  Final Research Proposal — CSIE 116 (HCI / HAX)
%
%  Engine: XeLaTeX (for CJK author names via PingFang TC on macOS).
%  Build:  latexmk -xelatex main.tex      (or see README.md)
% ============================================================
\documentclass[11pt]{article}

\usepackage{fontspec}
\usepackage{geometry}
\geometry{a4paper, margin=1in}

% --- CJK support for Chinese author names (XeLaTeX) ---
% If xeCJK / the font is unavailable, comment the next two lines out
% and the document still compiles (author names fall back to romanized).
\usepackage{xeCJK}
\setCJKmainfont{PingFang TC}

\usepackage{microtype}
\usepackage{graphicx}
\usepackage{amssymb}
\usepackage{booktabs}
\usepackage{tabularx}
\usepackage{array}
\usepackage[table]{xcolor}
\usepackage{enumitem}
\usepackage[numbers,sort&compress]{natbib}
\usepackage[colorlinks=true,linkcolor=accent,citecolor=accent,urlcolor=accent]{hyperref}

\definecolor{accent}{HTML}{4A5A78}
\definecolor{rust}{HTML}{A6505A}
\definecolor{rowgray}{HTML}{F4F1EC}

% Colored section headings without titlesec (minimal TeX install).
\makeatletter
\renewcommand\section{\@startsection{section}{1}{\z@}%
  {-3.5ex \@plus -1ex \@minus -.2ex}{2.3ex \@plus .2ex}%
  {\normalfont\Large\bfseries\color{accent}}}
\renewcommand\subsection{\@startsection{subsection}{2}{\z@}%
  {-3.25ex \@plus -1ex \@minus -.2ex}{1.5ex \@plus .2ex}%
  {\normalfont\large\bfseries\color{accent}}}
\makeatother
\setlist{itemsep=2pt, topsep=3pt}
\renewcommand{\arraystretch}{1.25}
\newcolumntype{Y}{>{\raggedright\arraybackslash}X}

% Property-matrix marks.
\newcommand{\cmark}{\textcolor{accent}{\checkmark}}
\newcommand{\xmark}{\textcolor{rust}{\ensuremath{\times}}}
\newcommand{\pmark}{\textcolor{rust}{\ensuremath{\sim}}}

\title{\vspace{-1.2cm}\textbf{\color{accent}Parting Is a Sweet Sorrow}\\[0.35em]
{\large\color{rust} An Ambient Co-Presence Space for Friendships\\ Drifting Apart in Young Adulthood}\\[0.6em]
{\normalsize\itshape Final Research Proposal — CSIE 116 (HCI / HAX)}}

\author{白昕宸 \quad 蘇楷峻 \quad 顏孜芸\\[0.2em]
{\small Department of Computer Science and Information Engineering}}
\date{June 2026}

\begin{document}
\maketitle

\begin{abstract}
\noindent
Adult friendships rarely end in a fight. Far more often they \emph{drift}: as
classes split, people relocate, and life paths diverge, the daily texture of
incidental contact quietly disappears and a once-close bond fades without any
explicit breakup. This pattern --- \textbf{passive friendship dissolution} ---
produces a distinctive loneliness that is invisible to the people experiencing
it and to the tools meant to help them. We argue that current technologies make
the problem worse: social platforms run on a logic of \emph{broadcast,
performance, and comparison}, AI companions offer comfort but never the feeling
of \emph{being needed}, and friend-matching apps optimize for making
\emph{new} ties rather than \emph{maintaining} existing close ones. Across the
HCI literature, three design lineages --- ambient awareness, digital third
places, and asynchronous communication --- each solve one half of the problem
and leave the same gap open. We propose \textbf{Parting Is a Sweet Sorrow}, an
ambient co-presence space in which a small group of friends inhabit a shared,
projected ``night sky'' that they build together from privately meaningful
objects and a shared vocabulary of light, color, and sound. The system is
designed to let members feel the group is \emph{still there} without demanding
performance or reply. We pose three research questions --- on whether projected
shared space sustains felt co-presence (RQ1), on whether non-performative
ambient presence lowers comparison and performance pressure while preserving
connectedness (RQ2), and on whether the space changes the willingness to
maintain a drifting friendship (RQ3) --- and outline a mixed-methods evaluation
anchored on the Affective Benefits and Costs Questionnaire (ABC-Q).
\end{abstract}

% ============================================================
\section{Research Motivation}
% ============================================================

We frame the motivation with the \textbf{ABT} structure (\emph{And, But,
Therefore}) required by the project rubric.

\subsection{And --- friendships are load-bearing infrastructure of young adulthood}

Close friendships are not a luxury of young adulthood; they are the
developmental scaffolding on which it is built. Emerging adulthood (roughly
18--30) is the life stage of maximal identity exploration and instability
\citep{arnett2000}, and it is also a peak of the lifespan loneliness curve
\citep{qualter2015, barreto2021}. During exactly this window, friendships do the
heavy lifting: they are where young adults build the intimacy skills that
underwrite \emph{every} later relationship \citep{rawlins2017}, and the quality
of close friendship in adolescence remains a significant long-term predictor of
adult emotional health \citep{allen2024}. Mutual close friendships also act as a
\emph{buffer}, protecting socially anxious young people against concurrent
psychosocial risk \citep{erath2010}. In short, the bonds formed in this period
are the foundation for lifelong relational capacity.

\subsection{But --- friendships dissolve passively, and the drift is invisible}

\textbf{But} adult friendships are unusually fragile, and they break in a
particular way. \emph{Passive friendship dissolution} is the gradual,
unstructured ending of a friendship \emph{without confrontation or any explicit
breakup conversation} \citep{vieth2022}. The Adult Friendship Dissolution
Process Model holds that passive dissolution is steered by a combination of the
actor's personal traits (attachment style, personality), situational pressures
(distance, competing life demands), and perceptions of the friend's likely
response --- which together push people toward \emph{gradually reducing contact}
rather than openly ending the tie \citep{vieth2022}. A complementary
features\,$\rightarrow$\,processes\,$\rightarrow$\,outcomes model of friendship
shows that ``gradual drift'' is a distinct dissolution outcome, separate from
direct rejection \citep{santucci2025}. Relocation studies confirm the mechanism
empirically: most adolescents experience a measurable period of diminished
access to companionship and intimacy after a move, with the effect larger for
those already socially anxious \citep{vernberg2006}.

The cruelty of this pattern is that it is \emph{invisible}. Because drift causes
no acute crisis, it never triggers help-seeking; by the time loneliness arrives,
it triggers \emph{withdrawal} rather than outreach \citep{cacioppo2015}. The
loneliness itself is a \emph{discrepancy} between the closeness one has and the
closeness one wants \citep{grigoriadis2024, hall2025} --- not an absence of
contact. People in constant digital contact with friends they love can
nonetheless feel the bond has quietly downgraded --- indeed, young adults
themselves report that online contact often lacks the authenticity of in-person
interaction and can enable avoidance and curation rather than genuine closeness
\citep{scott2022}. And the scale is large: a
majority of Gen~Z report feeling lonely ``sometimes or always''
\citep{cigna2020}, yet ``friendship drift'' remains essentially unaddressed.
Worse, the tools people reach for make it worse: a nine-year longitudinal study
found social-media use \emph{predicts increasing} loneliness over time, with the
strongest effect among those who use it specifically to maintain close ties
\citep{robertsyoungdavid2025}.

\subsection{Therefore --- make drift visible and restore ambient co-presence}

\textbf{Therefore} we propose to design for this specific loneliness: not the
loneliness of having no one, but the loneliness of friendships that have quietly
downgraded across diverging paths. Researchers have mapped \emph{who} is at risk
and \emph{what} the outcomes look like, but have not yet identified --- or
designed for --- the \emph{daily interaction patterns} that produce passive
drift \citep{santucci2025, hansen2025}. Our goal is an interactive space that
makes drift visible before it becomes irreversible and that restores a sense of
\emph{being together} without requiring the performance that today's tools
demand.

\subsection{Target audience}

Young adults aged roughly 18--30 (``adolescents in transition'') who are passing
through the life transitions --- graduation, relocation, class- or
school-changes, the start of work --- that scatter an existing friend group into
divergent daily routines. We target groups of close friends (group size $\sim$5)
who already share history and want to \emph{maintain} that closeness across
distance, rather than strangers seeking to form new ties.

% ============================================================
\section{Background: How Friendships Work and Why They Drift}
% ============================================================

To see why passive dissolution is hard to address, we locate it inside a model of
how friendship operates. \citet{santucci2025} describe friendship as a chain from
\emph{features} (individual behavior and cognition), through \emph{processes}
(positive interactions such as companionship and support; negative ones such as
conflict), to \emph{outcomes} (maintenance vs.\ dissolution, where dissolution
splits into \emph{gradual drift} and \emph{direct rejection}). Passive
dissolution is the gradual-drift branch: it is produced not by any negative
interaction but by the slow \emph{subtraction} of positive ones.

This is what makes it hard to design for. The field has mapped the
\emph{endpoints} --- who is at risk (young adults in transition: relocation,
graduation, school change \citep{vieth2022, vernberg2006}) and what the final
outcome looks like (a friendship observed only after it has dissolved) --- but
the \emph{middle}, the week-to-week interaction dynamics that tip a bond from
``maintained'' into ``drifted,'' remains largely unobserved \citep{santucci2025}.
\textbf{It is hard to design interventions for a process we cannot see.}

\paragraph{Why this loneliness matters --- a three-horizon view.}
The stakes compound over time. \textbf{Now}, the gap between attained and desired
closeness produces loneliness and depressed mood \citep{grigoriadis2024}.
\textbf{In the mid-term}, the person loses the buffer that good mutual
friendships provide against social anxiety \citep{erath2010}. \textbf{In the
long term}, because adolescent and young-adult friendship quality predicts adult
emotional health \citep{allen2024}, the damage carries forward. If young adults
internalize the lesson that ``relationships naturally fade,'' they may enter
adulthood unwilling to invest in any bond that requires sustained effort --- a
mass-scale system failure that each individual experiences as personal failure.

% ============================================================
\section{Related Work}
% ============================================================

\subsection{Loneliness as discrepancy, not absence}

A consistent finding across the loneliness literature is that loneliness is a
\emph{perceived discrepancy} between desired and actual closeness, not a count
of contacts \citep{grigoriadis2024, qualter2015}. Emerging adulthood is both the
developmental period of greatest relational flux \citep{arnett2000} and a peak in
the lifespan loneliness curve \citep{barreto2021, qualter2015}. \citet{hall2025}
characterize young-adult loneliness as a product of ``rapid life changes and a
lack of relational permanence and routine'' rather than of having no one ---
which is exactly the experiential signature of passive drift. This reframes the
design target: the problem is not connecting strangers but \emph{sustaining felt
closeness} in existing bonds whose routines have been disrupted.

\subsection{Passive vs.\ active dissolution}

\citet{vieth2022} distinguish \emph{active} dissolution (an explicit, often
conflict-driven ending) from \emph{passive} dissolution (a gradual withdrawal of
contact), and model the actor-, partner-, and context-level predictors that
route a friendship toward one or the other. \citet{santucci2025} situate
dissolution within the broader friendship lifecycle and identify gradual drift
as a distinct outcome with its own dynamics. Recent HCI work corroborates this
from the lived-experience side: interviews with people in long-term online
friendships find that such ties are sustained through accumulated, low-key
interaction and quietly erode when that interaction lapses, suggesting that
systems should support \emph{sustained interaction over time} rather than
momentary tie-strength \citep{yong2025}. Together these works establish that
passive dissolution is a phenomenon in its own right --- and that it has been
characterized at the level of \emph{predictors} and \emph{outcomes} but not yet
at the level of \emph{designable daily interaction}.

\subsection{Why current tools fail}

Existing technologies fail this problem in four characteristic ways.

\begin{description}[leftmargin=1.2em, style=nextline]
\item[The ``new friends'' paradox.] Friend-matching platforms optimize the
\emph{match} but ignore the roughly 200 hours of shared time needed to build a
close bond \citep{hall2019}; they manufacture introductions, not maintenance.
\item[Social snacks, not nutrition.] Mainstream social media runs on broadcast,
performance, and comparison. It fosters rumination about how one is perceived and
unfavorable social comparison \citep{stritzke2004}, and longitudinally
\emph{increases} loneliness for the very people using it to stay close
\citep{robertsyoungdavid2025}. It offers ``social snacks'' where the need is
emotional nutrition.
\item[Comfort without being needed.] AI companions can provide support and
comfort \citep{ta2020}, but cannot offer the feeling of \emph{being needed} by
another person --- the reciprocal interdependence that anchors a genuine close
bond.
\item[Breadth over depth.] Interventions and platforms skew toward
\emph{making new} ties or treating clinical loneliness, and toward breadth
metrics (follower counts, group-chat size). Tools for \emph{maintaining} existing
close bonds across distance, and for making drift visible before it is
irreversible, are largely missing \citep{akhterkhan2020, hansen2025}.
\end{description}

\subsection{Three HCI lineages, each solving half the problem}

The raw materials for a solution already exist in HCI, distributed across three
lineages --- but each addresses only one half of passive drift, and each leaves
the same gap (Table~\ref{tab:lineages}).

\begin{table}[t]
\centering
\caption{Three design lineages that each solve half of the passive-drift problem.}
\label{tab:lineages}
\small
\begin{tabularx}{\textwidth}{>{\bfseries}p{2.9cm} Y Y}
\toprule
Lineage & What it contributes & The gap it leaves \\
\midrule
Ambient awareness &
Peripheral cues make users aware of another's ongoing presence \emph{without
interrupting} them, raising felt connectedness \citep{deyguzman2006,
markopoulos2004, madianou2016}. &
Built almost entirely for \emph{one-to-one} or family ties (caregiving, couples),
not for peer \emph{groups} drifting apart. \\
\rowcolor{rowgray}
Digital third places &
Persistent shared ``places'' sustain repeated, low-pressure interaction and
casual co-presence \citep{kim2025}. &
Organized around outward \emph{interests} and often strangers, not around the
shared history of an existing friend group. \\
Asynchronous communication &
Removing the demand for real-time response lowers social anxiety and impression-
management load \citep{high2009, stritzke2004}. &
Treated as a \emph{feature} of messaging tools, never integrated into the design
of a shared interactive \emph{space}. \\
\bottomrule
\end{tabularx}
\end{table}

\subsection{The research gap}

No existing system unifies all four properties that passive drift demands:
\emph{(i)} ambient, peripheral presence that does not interrupt; \emph{(ii)}
\emph{group}-scale rather than dyadic; \emph{(iii)} asynchronous and
\emph{non-performative}, removing comparison and reply pressure; and
\emph{(iv)} organized around an existing \emph{relationship} and its shared
history rather than around interests or strangers.

Recent work has begun to move toward this gap without closing it.
\citet{shakeri2024} formalize \emph{passive co-presence} as a design construct
--- a low-effort, non-demanding sense of another's presence over distance ---
but study families rather than peer groups. \citet{leong2023} field-test an
ecosystem of distributed, physically-grounded artifacts for staying lightly
connected to remote friends, showing real demand for ambient, non-screen-centric
connection, yet realize it as discrete \emph{objects} rather than a single shared
\emph{space} that a group builds and inhabits together. And \citet{huschke2025}
distill design themes specifically for long-distance \emph{friendship}
relatedness, confirming that the friendship case --- as opposed to romantic or
family ties --- is a distinct, under-served target. The design space of ambient
co-presence has thus been developed productively for transnational families
\citep{madianou2016} and is only beginning to be charted for friends; no system
yet unifies the four properties above for a drifting peer group. We propose to
extend \emph{passive co-presence} \citep{shakeri2024} into exactly this gap.

% ============================================================
\section{Design Concept: \textit{Parting Is a Sweet Sorrow}}
% ============================================================

\subsection{An ambient co-presence space}

We propose a shared, persistent, projected space that a small group of friends
\emph{builds together}. The guiding metaphor is a \textbf{tent pitched with
friends and a shared night sky}. Inside the tent are only things the group
mutually understands --- a toy everyone once chipped in to buy, a pen someone
always used --- objects that carry shared memory. Step outside the tent and the
sky overhead is full of stars, where \emph{each star is a small spark sent by a
friend}. The space is, in one line, \emph{``an ambient atmosphere device, made
of the objects that hold your shared memories, that works across distance.''}

Crucially, the canvas is legible \emph{only to the group}. Following the
night-sky example: the group shares one sky-canvas; each member can add to it;
and what everyone finally sees is the same sky --- \emph{``the sky you built
together.''} Presence is felt in the periphery, not consumed in the center
\citep{weiserbrown1996}, and is expressed through graspable and ambient media
rather than text \citep{ishiiullmer1997}. The design directly inherits the
``presence display'' tradition that converts \emph{awareness} into
\emph{connectedness} \citep{deyguzman2006, markopoulos2004}.

\subsection{Three mechanisms}

\begin{enumerate}[label=\textbf{M\arabic*.}, leftmargin=2.4em]
\item \textbf{A shared interaction vocabulary.} The group agrees on what its
signals \emph{mean} --- a light or color, a type of object, an ambient sound.
Meaning is co-defined and local to the group, not dictated by the platform.
\item \textbf{Private hints.} Signals carry semantics private to the group ---
e.g., \emph{a purple shimmer might mean ``thinking of you.''} An outsider sees
only atmosphere; insiders read meaning. This is privacy by \emph{legibility}, not
by settings.
\item \textbf{Add / remove objects --- trace design.} Members add and remove
objects in the space; what accumulates is a \emph{trace} of presence over time,
a slow-built record of having-been-here-together rather than a feed of
performances.
\end{enumerate}

These mechanisms instantiate the four gap-closing properties: they are ambient
(M1, M2), group-scale and relationship-anchored (the shared sky, M3),
asynchronous and non-performative (traces accrue without demanding reply), and
they encode privacy architecturally through shared-only legibility (M2).

\subsection{Positioning against existing systems}

Table~\ref{tab:landscape} situates the concept against representative systems in
the surrounding design landscape. Each prior system supplies one ingredient;
none combines all four.

\begin{table}[t]
\centering
\caption{Design landscape as a property matrix. Columns are the four properties passive drift demands:
\textbf{Amb}ient (peripheral, non-interrupting), \textbf{Grp} (group- not dyadic-scale),
\textbf{A/NP} (asynchronous and non-performative), \textbf{Rel}ationship-anchored (an existing bond,
not interests or strangers). \cmark{}~satisfies, \pmark{}~partial, \xmark{}~absent. Each prior system
misses at least one; only our system satisfies all four.}
\label{tab:landscape}
\small
\setlength{\tabcolsep}{4pt}
\begin{tabularx}{\textwidth}{>{\raggedright\arraybackslash}p{4.2cm} c c c c >{\raggedright\arraybackslash}X}
\toprule
Representative system & Amb & Grp & A/NP & Rel & Key limitation for passive drift \\
\midrule
Digital Family Portrait \citep{mynatt2001} & \cmark & \xmark & \cmark & \cmark & One-to-one / caregiving, not a peer group. \\
\rowcolor{rowgray}
Family awareness, e.g.\ ASTRA \citep{markopoulos2004} & \cmark & \pmark & \cmark & \cmark & Built and validated for \emph{families}, not peers. \\
BuddyBeads \citep{kikingil2006} & \pmark & \xmark & \cmark & \cmark & Dyadic; fixed small clique, narrow expressivity. \\
\rowcolor{rowgray}
\emph{Animal Crossing} shared island & \xmark & \cmark & \cmark & \xmark & Game/activity-framed, not relationship maintenance. \\
\emph{Between} (couples) & \xmark & \xmark & \xmark & \cmark & Dyadic; creates response-waiting pressure. \\
\rowcolor{rowgray}
\emph{Discord} friend server & \xmark & \cmark & \pmark & \xmark & Performative posting, public comparison, topic-organized. \\
\midrule
\textbf{\textit{Parting Is a Sweet Sorrow}} & \cmark & \cmark & \cmark & \cmark & \textbf{Unifies all four for drifting peer groups.} \\
\bottomrule
\end{tabularx}
\end{table}

\subsection{Design principles}

The concept is governed by four principles drawn from the literature:
(1)~\emph{Periphery over center} --- presence lives at the edge of attention, not
as a notification demanding the center \citep{weiserbrown1996};
(2)~\emph{Ambient and tangible over textual} --- meaning is carried by light,
color, sound, and objects \citep{ishiiullmer1997};
(3)~\emph{Awareness becomes connectedness} --- the design target is felt
connectedness, the documented effect of well-made presence displays
\citep{deyguzman2006, markopoulos2004}; and
(4)~\emph{Non-performative by construction} --- because the medium cannot be used
to broadcast or compete, it removes the impression-management load that drives
both social anxiety \citep{high2009, stritzke2004} and the loneliness paradox of
mainstream social media \citep{robertsyoungdavid2025}. That such load can be
\emph{designed away} is not merely hypothetical: BeReal's design --- randomized
timing, reciprocal posting, and no curation --- has been shown to suppress
self-presentation pressure and reframe social comparison among young adults
\citep{kim2024bereal}, a precedent our atmosphere-from-objects mechanism extends
from momentary snapshots to persistent ambient presence.

% ============================================================
\section{Research Objectives and Questions}
% ============================================================

\noindent\textbf{Objective.} To design and evaluate an ambient, non-performative
shared space that lets a scattered group of close friends sustain felt
co-presence across distance, and to test whether doing so lowers performance and
comparison pressure while preserving --- or strengthening --- the willingness to
maintain a drifting friendship.

\begin{description}[leftmargin=1.4em, style=nextline]
\item[RQ1 (Foundation).] Can a projected shared space let members
\emph{persistently perceive that ``the friend group still exists''} without
requiring direct, real-time interaction?
\item[RQ2 (Core).] Can this blurry, non-performative co-presence \emph{maintain
connectedness while lowering performance and social-comparison pressure}, relative
to existing tools?
\item[RQ3 (Real world).] How does living with this space \emph{affect young
adults' willingness and intention to maintain a drifting friendship}, and through
what experienced mechanism?
\end{description}

The three questions are deliberately layered: RQ1 tests the \emph{precondition}
(felt persistent presence), RQ2 tests the \emph{core claim} (presence without
pressure), and RQ3 tests the \emph{downstream relational consequence} in lived
use.

% ============================================================
\section{Evaluation Plan}
% ============================================================

We propose a mixed-methods evaluation, with one strand per research question
(Table~\ref{tab:eval}).

\begin{table}[t]
\centering
\caption{Evaluation plan, by research question.}
\label{tab:eval}
\small
\begin{tabularx}{\textwidth}{>{\bfseries}p{1.4cm} Y p{4.2cm}}
\toprule
RQ & Approach & Primary measure \\
\midrule
RQ1 & Test whether the tent atmosphere and shared sky produce persistent awareness that ``the group is always there.'' We measure the \emph{sense of continuous presence}, not the \emph{number} of signals received. & Judgment-accuracy of perceived presence; self-reported continuity of awareness. \\
\rowcolor{rowgray}
RQ2 & Compare our space against participants' existing tools. Hypothesis: \emph{affective benefits are maintained while affective costs drop.} & Affective Benefits and Costs Questionnaire (ABC-Q): Affective Benefits, Affective Costs, Perceived Awareness \citep{ijsselsteijn2009}. \\
RQ3 & Qualitative interviews: does a presence that requires \emph{no} image-management feel more relaxed, and does it sustain connection? How does this happen? & Thematic analysis of interview transcripts. \\
\bottomrule
\end{tabularx}
\end{table}

The ABC-Q \citep{ijsselsteijn2009} is central to RQ2 because it was purpose-built
to measure the \emph{affective} characteristics of mediated awareness --- both
its benefits (e.g., connectedness) and its costs (e.g., obligation, pressure) ---
which is exactly the benefit-without-cost trade-off our design claims to improve.
RQ1's emphasis on \emph{accuracy of perceived presence} rather than message
volume operationalizes the calm-technology claim that the system should inform
the periphery without capturing the center \citep{weiserbrown1996}. RQ3's
qualitative strand is where we expect to surface the experienced mechanism ---
whether removing performance is what makes sustained, low-effort closeness feel
possible.

% ============================================================
\section{Self-Assessment: McGee's Eight Questions}
% ============================================================

\begin{table}[t]
\centering
\caption{Self-assessment against Prof.\ Kevin McGee's eight fundamental questions.}
\label{tab:mcgee}
\small
\begin{tabularx}{\textwidth}{>{\bfseries}p{4.6cm} Y}
\toprule
Question & Our answer \\
\midrule
What is the research problem? & How to sustain felt co-presence in friend groups undergoing \emph{passive} dissolution, without the performance and comparison pressure of current tools. \\
\rowcolor{rowgray}
Is the problem significant? & Yes: young-adult loneliness is a peak-vulnerability, high-prevalence phenomenon \citep{barreto2021, cigna2020} with lifelong consequences for relational capacity \citep{allen2024, rawlins2017}. \\
What is the contribution to knowledge? & A design and an evaluation of \emph{group-scale, non-performative ambient co-presence} for drifting peer friendships --- a gap left by all three existing HCI lineages. \\
\rowcolor{rowgray}
Is the contribution vital? & Drift is invisible and currently un-designed-for; existing tools \emph{worsen} it \citep{robertsyoungdavid2025}. A maintenance-oriented, non-performative alternative is needed. \\
Originality of the contribution. & Extends the \emph{passive co-presence} construct \citep{shakeri2024} and ambient co-presence \citep{madianou2016} from families to peer groups, and unifies ambient $+$ group $+$ asynchronous $+$ relationship-anchored design --- a combination not present in prior systems (Table~\ref{tab:landscape}). \\
\rowcolor{rowgray}
How do readers know it is original? & The landscape and lineage analyses (Tables~\ref{tab:lineages},~\ref{tab:landscape}) show each prior system supplies only one ingredient; none combines all four. \\
Validity of the contribution. & To be established through a layered, falsifiable design (RQ1--RQ3) using a validated instrument (ABC-Q, \citealp{ijsselsteijn2009}) plus thematic analysis. \\
\rowcolor{rowgray}
Is it appropriate for the discipline? & Yes: it is squarely an HCI / HAX contribution, building on calm technology \citep{weiserbrown1996}, tangible bits \citep{ishiiullmer1997}, and presence displays \citep{deyguzman2006}. \\
\bottomrule
\end{tabularx}
\end{table}

% ============================================================
\section{Expected Contributions}
% ============================================================

We expect three contributions. \emph{Conceptually}, we name and target a
specific, under-served form of loneliness --- the loneliness of passive
friendship drift --- and articulate why mainstream tools structurally worsen it.
\emph{In design}, we contribute \textit{Parting Is a Sweet Sorrow}: a
group-scale, non-performative ambient co-presence space defined by three
mechanisms (shared vocabulary, private hints, trace-based objects) that together
close the gap left by prior ambient-awareness, third-place, and asynchronous
systems. \emph{Methodologically}, we contribute an evaluation design that
separates felt persistent presence (RQ1) from the benefit-without-cost claim
(RQ2, via ABC-Q) and from the downstream relational effect (RQ3), giving each a
distinct, testable measure.

\vspace{0.4em}
\noindent\textit{Parting is such sweet sorrow} --- the aim of this work is to
make the sorrow gentler, and the parting less final, by letting friends remain
ambiently, unperformingly, together.

% ============================================================
\bibliographystyle{plainnat}
\bibliography{references}
% ============================================================

\end{document}
